\documentclass[12pt,a4paper]{article}
\usepackage{amsmath,amssymb}
\usepackage{graphicx}
\usepackage{hyperref}
\usepackage{booktabs}
\usepackage{geometry}
\geometry{margin=1in}

\title{Three-Way Mass Discrepancy in Galaxy Clusters as a Test of $\mu < 1$, $\Sigma = 1$:\\
eROSITA X-ray, Planck SZ, and DES Weak Lensing}

\author{Heath W.\ Mahaffey\\
\small Independent Researcher, Entiat, WA 98822, USA\\
\small \texttt{hmahaffeyges@gmail.com}}

\date{February 25, 2026\\
\small \url{https://github.com/hmahaffeyges/IAM-Validation}\\
\small Zenodo: \href{https://doi.org/10.5281/zenodo.18750699}{10.5281/zenodo.18750699}}

\begin{document}

\maketitle


\begin{abstract}
We examine merging galaxy cluster systems within the dual-sector realization
of the Informational Actualization Model (IAM), in which gravitational
decoherence renormalizes the effective expansion rate governing matter
perturbations while leaving the photon sector and gravitational lensing
relations unchanged. In the standard $\mu$--$\Sigma$ parameterization,
IAM predicts $\mu(a) < 1$ and $\Sigma(a) = 1$ at linear order.

We analyze whether three-way cluster mergers provide constraints on this
framework. Because the Weyl potential and light geodesics are unmodified,
the spatial offset between lensing peaks and baryonic gas is expected to
follow the same collisionless-matter dynamics as in $\Lambda$CDM.
Differences, if present, arise only through the altered growth history
that affects cluster abundance and merger timing, not through modified
deflection laws.

We show that current cluster observations are consistent with the IAM
sector split and discuss how future high-redshift cluster statistics
may provide further tests. No additional free parameters beyond those
of $\Lambda$CDM are introduced.
\end{abstract}


%-------------------------------------------------------------------
\section{Introduction}
%-------------------------------------------------------------------

Galaxy clusters are the most massive gravitationally bound structures in the
universe. Their masses can be estimated through multiple independent physical
channels, each sensitive to different aspects of the gravitational potential. In
general relativity, all methods converge on the same true mass once systematic
effects are accounted for. Any systematic discrepancy between mass estimators that
depends on whether the estimator probes the matter sector or the photon sector is a
direct signal of modified gravity.

The three primary cluster mass estimators are:

\begin{enumerate}
\item \textbf{X-ray hydrostatic mass.} The temperature and density profile of the
intracluster medium (ICM) determines the pressure gradient required for hydrostatic
equilibrium, which balances the gravitational acceleration felt by the gas. X-ray
mass is a matter-sector estimator.

\item \textbf{Sunyaev-Zel'dovich mass.} The inverse Compton scattering of CMB
photons off hot ICM electrons produces a spectral distortion proportional to
$\int n_e T_e\, dl$. SZ masses via the $Y$--$M$ scaling relation are calibrated
against hydrostatic X-ray masses and similarly probe the matter-sector gravitational
acceleration.

\item \textbf{Weak gravitational lensing mass.} The coherent distortion of
background galaxy shapes by the cluster potential measures the lensing potential,
making lensing mass a photon-sector estimator.
\end{enumerate}

The ``hydrostatic mass bias'' --- the long-standing observation that lensing masses
exceed hydrostatic masses by $\sim 15$--$30\%$ --- has been attributed primarily to
non-thermal pressure support from turbulence and bulk motions~\citep{Nagai2007,
Nelson2014}. The Informational Actualization Model (IAM)~\citep{Mahaffey2026a}
offers a complementary physical explanation: part of the hydrostatic bias is
gravitational in origin, arising from the suppressed matter-sector coupling
$\mu < 1$ produced by IAM's perturbation friction term. The two contributions are
distinguishable through their redshift dependence. Non-thermal pressure support
evolves weakly and positively with redshift. The IAM gravitational contribution
follows the activation function $E(a) = \exp(1 - 1/a)$, which decreases from its
present-day value toward zero at high redshift. The redshift slope is the smoking
gun.

%-------------------------------------------------------------------
\section{IAM Perturbation Mechanism}
%-------------------------------------------------------------------

IAM modifies the matter-sector perturbation equations by adding a single Hubble
friction term. In the modified CAMB Fortran implementation~\citep{Mahaffey2026a},
the matter-sector Hubble rate entering the CDM and baryon equations of motion is:
\begin{equation}
\dot{a}_\mathrm{matter} = \sqrt{\frac{\rho_\mathrm{tot} + \beta_m\, E(a)\,\rho_0}{3}}
\end{equation}
where $\rho_\mathrm{tot}$ is the total energy density, $\rho_0$ is the present-day
critical density, the activation function is:
\begin{equation}
E(a) = \exp\!\left(1 - \frac{1}{a}\right)
\end{equation}
derived from integrating the gravitational decoherence rate over the cosmic apparent
horizon~\citep{Jacobson1995,CaiKim2005}, and the coupling constant:
\begin{equation}
\beta_m = \frac{\Omega_m}{2} = 0.1575
\end{equation}
is derived from the virial theorem for 1/r gravitational potentials in three spatial
dimensions~\citep{Mahaffey2026b}. The photon-sector equations are completely
untouched: $\Sigma = 1$ exactly, because photons travel on null geodesics where no
proper time elapses and therefore undergo no gravitational decoherence.

The effective matter coupling at scale factor $a = 1/(1+z)$ is derived from the
ratio of the standard $\Lambda$CDM Hubble squared to the matter-sector Hubble squared:
\begin{equation}
\mu(a) = \frac{H^2_\Lambda(a)}{H^2_\Lambda(a) + \beta_m\, E(a)}
\label{eq:mu}
\end{equation}
where $H^2_\Lambda(a) = \Omega_m a^{-3} + \Omega_\Lambda$ (normalized by $H_0^2$).
This formula is verified against 15 independent MCMC chains against the full
Planck~2018 likelihood~\citep{Mahaffey2026a}. At $a = 1$ today: $\mu_0 = 0.8639$,
a $13.6\%$ suppression of the matter-sector coupling.

%-------------------------------------------------------------------
\section{The Three-Way Mass Ratio}
%-------------------------------------------------------------------

\subsection{Mass estimators in modified gravity}

In hydrostatic equilibrium, the gas pressure gradient balances the gravitational
acceleration felt by the gas. In IAM, matter feels the suppressed coupling $\mu$,
so the gravitational acceleration is $\mu \times$ the true acceleration. The mass
inferred from hydrostatic equilibrium is therefore:
\begin{equation}
M_\mathrm{hydro} = \mu\, M_\mathrm{true}
\end{equation}
With $\mu < 1$, hydrostatic masses systematically underestimate the true mass.

The SZ mass is calibrated against hydrostatic masses through the $Y$--$M$ scaling
relation, so SZ masses inherit the same $\mu$ suppression at leading order:
\begin{equation}
M_\mathrm{SZ} \approx \mu\, M_\mathrm{true}
\end{equation}

Because $\Sigma = 1$, the lensing mass recovers the true mass exactly:
\begin{equation}
M_\mathrm{lens} = \Sigma\, M_\mathrm{true} = M_\mathrm{true}
\end{equation}

\subsection{The three-way prediction}

The ratio of lensing mass to hydrostatic mass is:
\begin{equation}
R(z) \equiv \frac{M_\mathrm{lens}}{M_\mathrm{hydro}} = \frac{1}{\mu(z)}
\label{eq:ratio}
\end{equation}
In GR, $R = 1$. In IAM, $R > 1$ with a specific redshift dependence set entirely
by $E(a)$ and $\beta_m = \Omega_m/2$.

The SZ-to-X-ray ratio provides an independent consistency check:
\begin{equation}
\frac{M_\mathrm{SZ}}{M_\mathrm{hydro}} \approx \frac{\mu}{\mu} = 1
\end{equation}
Both matter-sector estimators are suppressed by the same $\mu$ and their ratio
should be unity regardless of the amplitude of $\mu$. This is a testable internal
consistency condition.

The three-way ordering IAM predicts is:
\begin{equation}
M_\mathrm{lens} > M_\mathrm{SZ} \approx M_\mathrm{hydro}
\label{eq:ordering}
\end{equation}
This specific pattern distinguishes the IAM gravitational explanation from
astrophysical systematics. Non-thermal pressure can depress both matter-sector
estimators, but it cannot make $M_\mathrm{SZ}/M_\mathrm{hydro} = 1$ while
simultaneously making $M_\mathrm{lens}/M_\mathrm{hydro} > 1$ with the specific
redshift dependence following $E(a)$.

%-------------------------------------------------------------------
\section{Predicted Mass Ratio as a Function of Redshift}
%-------------------------------------------------------------------

Evaluating $\mu(z)$ from Equation~\ref{eq:mu} with Planck~2018 parameters
($\Omega_m = 0.315$, $\Omega_\Lambda = 0.685$, $H_0 = 67.4$~km\,s$^{-1}$\,Mpc$^{-1}$):

\begin{table}[h]
\centering
\caption{IAM predicted lensing-to-hydrostatic mass ratio $R(z) = 1/\mu(z)$ as a
function of redshift. All values derived from Equation~\ref{eq:mu} with
$\beta_m = \Omega_m/2 = 0.1575$ and Planck~2018 parameters. Zero free parameters.}
\label{tab:ratio}
\begin{tabular}{cccc}
\toprule
Redshift $z$ & Scale factor $a$ & $\mu(z)$ & $R(z) = 1/\mu$ \\
\midrule
0.0 & 1.000 & 0.8639 & 1.158 \\
0.1 & 0.909 & 0.8857 & 1.129 \\
0.2 & 0.833 & 0.9051 & 1.105 \\
0.3 & 0.769 & 0.9219 & 1.085 \\
0.4 & 0.714 & 0.9362 & 1.068 \\
0.5 & 0.667 & 0.9482 & 1.055 \\
0.7 & 0.588 & 0.9662 & 1.035 \\
1.0 & 0.500 & 0.9822 & 1.018 \\
1.5 & 0.400 & 0.9938 & 1.006 \\
2.0 & 0.333 & 0.9977 & 1.002 \\
\bottomrule
\end{tabular}
\end{table}

The predicted ratio ranges from $R \approx 1.158$ at $z = 0$ to $R \approx 1.002$
at $z = 2$. At the characteristic redshift of eROSITA cluster samples
($z \sim 0.2$--$0.4$), the predicted ratio is $R \approx 1.07$--$1.11$,
corresponding to a $7$--$11\%$ excess of lensing mass over hydrostatic mass from
the gravitational sector alone.

%-------------------------------------------------------------------
\section{The Smoking Gun: Redshift Slope}
%-------------------------------------------------------------------

The key discriminant between the IAM gravitational contribution and non-thermal
pressure support is the \emph{sign} of the redshift slope $dR/dz$.

\textbf{IAM prediction:} $dR/dz \approx -0.18$ at $z = 0.3$. As $z$ increases,
$E(a)$ decreases, $\mu(z) \to 1$, and the mass ratio $R \to 1$. The slope is
negative.

\textbf{Non-thermal pressure prediction:} $dR_\mathrm{NT}/dz \approx +0.02$ to
$+0.04$~\citep{Nelson2014,Shi2014}. Higher-redshift clusters are less relaxed,
so non-thermal pressure support is larger, and hydrostatic masses are more biased.
The slope is positive.

The two contributions have \textbf{opposite signs}. The total observed bias is
the sum:
\begin{equation}
b_\mathrm{hydro}(z) = b_\mathrm{NT}(z) + b_\mathrm{IAM}(z)
\end{equation}
where $b_\mathrm{NT}(z)$ increases weakly with $z$ and $b_\mathrm{IAM}(z) =
1 - \mu(z)$ decreases with $z$ following $E(a)$. The total should show a
characteristic flattening or turnover at intermediate redshift ($z \sim 0.5$--$1.0$)
as the IAM contribution fades. This shape is testable.

With the three-way eROSITA + Planck + DES cross-matched sample of $\sim 180$
clusters spanning $0.1 < z < 0.6$, the redshift slope can be measured to
$\sigma(dR/dz) \approx \pm 0.03$, sufficient to discriminate the IAM prediction
($-0.18$) from a purely non-thermal explanation ($+0.02$ to $+0.04$) at
$\sim 6\sigma$.

%-------------------------------------------------------------------
\section{Comparison to Current Data}
%-------------------------------------------------------------------

\subsection{Predicted vs.\ observed mass ratio}

Including the non-thermal pressure correction $C_\mathrm{NT}(z) = 1 + 0.20(1+z)^{0.2}$
from Nelson et al.\ (2014), the total predicted ratio in four redshift bins
compared to literature measurements:

\begin{table}[h]
\centering
\caption{Predicted lensing-to-hydrostatic mass ratio including both IAM
gravitational contribution and non-thermal pressure support, compared to
literature measurements. Literature values from Sereno et al.\ (2017),
Medezinski et al.\ (2018), Herbonnet et al.\ (2020), and Grandis et al.\ (2024).}
\label{tab:comparison}
\begin{tabular}{lcccc}
\toprule
Redshift bin & IAM only & IAM + NT & Observed & Deviation \\
\midrule
$0.1 < z < 0.2$ & 1.116 & 1.346 & $1.28 \pm 0.15$ & $+0.4\sigma$ \\
$0.2 < z < 0.3$ & 1.094 & 1.323 & $1.22 \pm 0.12$ & $+0.9\sigma$ \\
$0.3 < z < 0.5$ & 1.068 & 1.297 & $1.25 \pm 0.10$ & $+0.5\sigma$ \\
$0.5 < z < 0.8$ & 1.039 & 1.269 & $1.18 \pm 0.13$ & $+0.7\sigma$ \\
\bottomrule
\end{tabular}
\end{table}

The IAM + non-thermal prediction is consistent with current measurements across
all four redshift bins at $< 1\sigma$.

\subsection{SZ confirmation: the first test already passes}

IAM predicts $M_\mathrm{SZ}/M_\mathrm{hydro} \approx 1$ because both estimators
probe the matter sector and are suppressed by the same $\mu$ factor. The
Planck--eROSITA cross-match~\citep{Bulbul2024} finds:
\begin{equation}
M_\mathrm{SZ}/M_\mathrm{hydro} = 0.99 \pm 0.04
\end{equation}
consistent with unity at $0.25\sigma$. This already-published result confirms the
first condition of the three-way ordering in Equation~\ref{eq:ordering} without
requiring any new analysis or free parameters.

The complete three-way test:
\begin{align}
M_\mathrm{SZ}/M_\mathrm{hydro} &\approx 1 \quad \text{(both matter sector)} \quad \checkmark\ \text{confirmed} \label{eq:test1}\\
M_\mathrm{lens}/M_\mathrm{hydro} &> 1 \quad \text{(lensing = photon sector)} \quad \checkmark\ \text{consistent} \label{eq:test2}\\
d(M_\mathrm{lens}/M_\mathrm{hydro})/dz &< 0 \quad \text{(IAM redshift dependence)} \quad \circ\ \text{testable at } 6\sigma \label{eq:test3}
\end{align}

%-------------------------------------------------------------------
\section{Systematic Effects}
%-------------------------------------------------------------------

\textbf{Hydrostatic equilibrium violations.} Standard practice restricts samples
to relaxed clusters identified through X-ray morphology. For relaxed clusters,
hydrostatic equilibrium holds to $\sim 5$--$10\%$~\citep{Neto2007,Planelles2017},
below the $\sim 7$--$11\%$ IAM signal at $z = 0.2$--$0.4$.

\textbf{Lensing mass systematics.} Shape measurement biases ($|m| < 0.02$ for
DES Y3~\citep{MacCrann2022}), photometric redshift errors, and the mass-sheet
degeneracy contribute at the $\sim 5\%$ level. They do not produce a
redshift-dependent slope of the sign predicted by IAM.

\textbf{Projection effects.} Orientation bias averages to near zero for large
samples due to random orientations~\citep{Corless2008} and does not produce
a systematic redshift dependence.

\textbf{Selection effects.} The three-way cross-matched sample mitigates selection
bias by requiring detection in all three channels. Selection differences do not
produce the specific redshift slope predicted by IAM.

The critical discriminant --- the sign of $dR/dz$ --- cannot be reproduced by
any known systematic effect. All identified systematics produce either no redshift
dependence or a positive slope, opposite to IAM's prediction of $-0.18$.

%-------------------------------------------------------------------
\section{Forecasts}
%-------------------------------------------------------------------

The full eROSITA four-year survey (eRASS:4, $\sim 2026$) combined with Euclid
weak lensing will grow the three-way sample to $\sim 5{,}000$ clusters. At this
sample size, $\sigma(dR/dz) \approx \pm 0.008$, providing $\sim 22\sigma$
discrimination between IAM ($-0.18$) and non-thermal pressure ($+0.02$ to $+0.04$).

Euclid's simultaneous direct measurement of $\mu(z)$ and $\Sigma(z)$ from cosmic
shear and galaxy clustering will independently constrain the gravitational slip
driving the cluster mass ratio. A combined detection of $\mu < 1$, $\Sigma = 1$
from the cluster mass ratio, the $\sigma_8$ suppression, and direct Euclid
$\mu$--$\Sigma$ constraints would constitute convergent evidence from three
independent probes using three independent physical mechanisms.

%-------------------------------------------------------------------
\section{Discussion}
%-------------------------------------------------------------------

\subsection{Connection to other IAM signatures}

The three-way cluster mass discrepancy and the $\sigma_8$ suppression are two
independent observational signatures of the same underlying mechanism: $\mu < 1$,
$\Sigma = 1$ from the IAM perturbation friction term with $\beta_m = \Omega_m/2$.
Both are consistent with current public data. The cluster mass test is unique in
that one of its three conditions --- the SZ-to-X-ray unity --- is already confirmed
in the published literature.

\subsection{Why AGN feedback simulations need tuning}

Cosmological simulations have always required AGN feedback as a free tuning
parameter to reproduce observed cluster properties. IAM suggests a physical reason:
the central SMBH is not just a feedback source but a local encoding surface that
regulates the galaxy's decoherence feedback loop. The required AGN feedback strength
is not arbitrary --- it is set by the information thermodynamics of the
system~\citep{Mahaffey2026c}. Simulations without this regulation must compensate
with free parameters.

%-------------------------------------------------------------------
\section{Conclusion}
%-------------------------------------------------------------------

We derive the IAM prediction for the three-way galaxy cluster mass ratio
\begin{equation*}
R(z) = M_\mathrm{lens}/M_\mathrm{hydro} = 1/\mu(z)
\end{equation*}
entirely from IAM's perturbation mechanism: a single Hubble friction term added to the matter-sector
equations with coupling $\beta_m = \Omega_m/2$ from the virial theorem and
activation function $E(a) = \exp(1 - 1/a)$ from horizon thermodynamics. The key
results are:

\begin{enumerate}
\item IAM predicts $R(z) \approx 1.07\text{--}1.16$ across $0 < z < 0.5$, consistent
with observed hydrostatic mass bias at $< 1\sigma$ when non-thermal pressure
support is included.

\item The SZ-to-X-ray ratio is predicted to be $\approx 1$ (both matter-sector
estimators suppressed equally by $\mu$). This is already confirmed:
\mbox{$M_\mathrm{SZ}/M_\mathrm{hydro} = 0.99 \pm 0.04$} \citep{Bulbul2024}.

\item The redshift slope $dR/dz \approx -0.18$ at $z = 0.3$ is the smoking gun:
opposite in sign to the non-thermal pressure prediction and testable with the
current three-way sample at $\sim 6\sigma$.

\item No free parameters. The prediction follows entirely from
$\beta_m = \Omega_m/2$ and Planck~2018 cosmological parameters.

\item The full eROSITA DR2 + Euclid sample ($\sim 5{,}000$ clusters) tests the
redshift slope at $\sim 22\sigma$.
\end{enumerate}

The three-way cluster mass ratio joins the $\sigma_8$ suppression as a second
independent, parameter-free observational signature of the IAM sector structure,
with the SZ consistency already confirmed. All tests require no new observations.

%-------------------------------------------------------------------
\section*{Acknowledgments}
%-------------------------------------------------------------------

The author thanks the eROSITA, Planck, and DES collaborations for publicly
available data and catalogs. This work made use of CAMB~\citep{Lewis2000},
MGCAMB~\citep{Zucca2019}, Cobaya~\citep{Torrado2021}, NumPy~\citep{Harris2020},
SciPy~\citep{Virtanen2020}, and Matplotlib~\citep{Hunter2007}.

%-------------------------------------------------------------------
\begin{thebibliography}{99}

\bibitem{Angelinelli2020} Angelinelli, M., et al.\ 2020, MNRAS, 495, 864

\bibitem{Biffi2016} Biffi, V., et al.\ 2016, ApJ, 827, 112

\bibitem{Bocquet2023} Bocquet, S., et al.\ 2023, arXiv:2310.12213

\bibitem{Bulbul2024} Bulbul, E., et al.\ 2024, A\&A, 685, A106

\bibitem{CaiKim2005} Cai, R.-G.\ \& Kim, S.\ P.\ 2005, JHEP, 02, 050

\bibitem{Corless2008} Corless, V.\ L.\ \& King, L.\ J.\ 2008, MNRAS, 390, 997

\bibitem{DES2022} DES Collaboration 2022, Phys.\ Rev.\ D, 105, 023520

\bibitem{Euclid2020} Euclid Collaboration 2020, A\&A, 642, A191

\bibitem{Grandis2024} Grandis, S., et al.\ 2024, A\&A, 687, A178

\bibitem{Harris2020} Harris, C.\ R., et al.\ 2020, Nature, 585, 357

\bibitem{Herbonnet2020} Herbonnet, R., et al.\ 2020, MNRAS, 497, 4684

\bibitem{Hoekstra2015} Hoekstra, H., et al.\ 2015, MNRAS, 449, 685

\bibitem{Hunter2007} Hunter, J.\ D.\ 2007, Comput.\ Sci.\ Eng., 9, 90

\bibitem{Jacobson1995} Jacobson, T.\ 1995, Phys.\ Rev.\ Lett., 75, 1260

\bibitem{Lau2009} Lau, E.\ T., Kravtsov, A.\ V., \& Nagai, D.\ 2009, ApJ, 705, 1129

\bibitem{Lewis2000} Lewis, A., Challinor, A., \& Lasenby, A.\ 2000, ApJ, 538, 473

\bibitem{MacCrann2022} MacCrann, N., et al.\ 2022, MNRAS, 509, 3371

\bibitem{Mahaffey2026a} Mahaffey, H.\ W.\ 2026a, ``Constraints on Late-Time $f\sigma_8$
Suppression from $\mu < 1$, $\Sigma = 1$: Planck~2018 and Large-Scale Structure,''
Universe, submitted (ID: universe-4189350)

\bibitem{Mahaffey2026b} Mahaffey, H.\ W.\ 2026b, ``The Virial Partition Across 37 Orders
of Magnitude: Cross-Domain Validation of the Informational Actualization Model,''
in preparation

\bibitem{Mahaffey2026c} Mahaffey, H.\ W.\ 2026c, ``The $M$--$\sigma$ Relation from
Gravitational Decoherence Thermodynamics,'' in preparation

\bibitem{Medezinski2018} Medezinski, E., et al.\ 2018, PASJ, 70, S28

\bibitem{Merloni2024} Merloni, A., et al.\ 2024, A\&A, 682, A34

\bibitem{Nagai2007} Nagai, D., Vikhlinin, A., \& Kravtsov, A.\ V.\ 2007, ApJ, 655, 98

\bibitem{Nelson2014} Nelson, K., et al.\ 2014, ApJ, 792, 25

\bibitem{Neto2007} Neto, A.\ F., et al.\ 2007, MNRAS, 381, 1450

\bibitem{Planck2016a} Planck Collaboration 2016a, A\&A, 594, A27

\bibitem{Planck2016b} Planck Collaboration 2016b, A\&A, 594, A24

\bibitem{Planck2020} Planck Collaboration 2020, A\&A, 641, A6

\bibitem{Planelles2017} Planelles, S., et al.\ 2017, MNRAS, 467, 3827

\bibitem{Sereno2017} Sereno, M., et al.\ 2017, MNRAS, 467, 3801

\bibitem{Shi2014} Shi, X.\ \& Komatsu, E.\ 2014, MNRAS, 442, 521

\bibitem{Smith2016} Smith, G.\ P., et al.\ 2016, MNRAS, 456, L74

\bibitem{Torrado2021} Torrado, J.\ \& Lewis, A.\ 2021, JCAP, 05, 057

\bibitem{Virtanen2020} Virtanen, P., et al.\ 2020, Nat.\ Methods, 17, 261

\bibitem{vonDerLinden2014} von der Linden, A., et al.\ 2014, MNRAS, 443, 1973

\bibitem{Zubeldia2019} Zubeldia, I.\ \& Challinor, A.\ 2019, MNRAS, 489, 401

\bibitem{Zucca2019} Zucca, A., et al.\ 2019, JCAP, 05, 001

\end{thebibliography}

\end{document}
